% theme: applications_reading_numerical_tables
\MajorQuestionLesson{数表を読み取るような問題}
\SetRowSpace{5mm}

\TypeQuestion{表から数量の関係を読み取り、連立方程式をつくって問いに答えなさい。}{tokuten}

%シュートの名前自体に点数が入ってしまっていて、表を読む問題になってないです。
\QQTall{あるバスケットボールの試合で、2点シュートと3点シュートが合わせて20本成功し、それらによる得点は47点であった。成功した2点シュートと3点シュートは、それぞれ何本か。
\par\smallskip
\begin{center}
\begin{tabular}{|c|c|c|}\hline
 & 2点シュート & 3点シュート \\ \hline
1本あたりの得点 & $2$点 & $3$点 \\ \hline
\end{tabular}
\end{center}}{}

\QQTall{25問のクイズで、正解は1問につき4点、不正解は1問につき1点減点、無解答は0点である。無解答が4問で、得点が59点であったとき、正解と不正解はそれぞれ何問か。
\par\smallskip
\begin{center}
\begin{tabular}{|c|c|c|c|}\hline
 & 正解 & 不正解 & 無解答 \\ \hline
1問あたりの得点 & $4$点 & $-1$点 & $0$点 \\ \hline
\end{tabular}
\end{center}}{}

\QQTall{ある大会で14試合を行い、勝ち、引き分け、負けに対して表のように勝ち点が与えられる。3試合に負け、勝ち点の合計が25点であったとき、勝ちと引き分けはそれぞれ何試合か。
\par\smallskip
\begin{center}
\begin{tabular}{|c|c|c|c|}\hline
 & 勝ち & 引き分け & 負け \\ \hline
1試合あたりの勝ち点 & $3$点 & $1$点 & $0$点 \\ \hline
\end{tabular}
\end{center}}{}

\QQTall{クラスAとクラスBの生徒が合わせて50人いる。テストの平均点が表のとおりであるとき、各クラスの生徒数を求めなさい。
\par\smallskip
\begin{center}
\begin{tabular}{|c|c|c|c|}\hline
 & クラスA & クラスB & 全体 \\ \hline
平均点 & $65$点 & $80$点 & $71$点 \\ \hline
\end{tabular}
\end{center}}{}

\QQTall{種類Aと種類Bの荷物が合わせて30箱ある。1箱あたりの平均の重さが表のとおりであるとき、それぞれ何箱あるか。
\par\smallskip
\begin{center}
\begin{tabular}{|c|c|c|c|}\hline
 & 種類A & 種類B & 全体 \\ \hline
1箱あたりの平均 & $4$kg & $7$kg & $5$kg \\ \hline
\end{tabular}
\end{center}}{}

\QQTall{ある学年の男子と女子が合わせて40人いる。通学時間の平均が表のとおりであるとき、男子と女子の人数をそれぞれ求めなさい。
\par\smallskip
\begin{center}
\begin{tabular}{|c|c|c|c|}\hline
 & 男子 & 女子 & 学年全体 \\ \hline
通学時間の平均 & $18$分 & $30$分 & $22.5$分 \\ \hline
\end{tabular}
\end{center}}{}
\TypeQuestion{表に示された1個あたりの数量と全体の数量から、それぞれの個数を求めなさい。}{kumiawase}

\QQTall{製品Aと製品Bを作ったところ、材料を34kg、作業時間を41分使った。製品Aと製品Bをそれぞれ何個作ったか。
\par\smallskip
\begin{center}
\begin{tabular}{|c|c|c|}\hline
 & 製品A & 製品B \\ \hline
1個あたりの材料 & $3$kg & $2$kg \\ \hline
1個あたりの作業時間 & $2$分 & $5$分 \\ \hline
\end{tabular}
\end{center}}{}

\QQTall{ギフトセットPとギフトセットQを作るために、鉛筆を32本、消しゴムを22個使った。セットPとセットQをそれぞれ何組作ったか。
\par\smallskip
\begin{center}
\begin{tabular}{|c|c|c|}\hline
 & セットP & セットQ \\ \hline
1組あたりの鉛筆 & $2$本 & $5$本 \\ \hline
1組あたりの消しゴム & $3$個 & $1$個 \\ \hline
\end{tabular}
\end{center}}{}

\QQTall{弁当Aと弁当Bを作るために、おにぎりを27個、おかずを41個用意した。弁当Aと弁当Bをそれぞれ何個作ったか。
\par\smallskip
\begin{center}
\begin{tabular}{|c|c|c|}\hline
 & 弁当A & 弁当B \\ \hline
1個あたりのおにぎり & $2$個 & $1$個 \\ \hline
1個あたりのおかず & $1$個 & $3$個 \\ \hline
\end{tabular}
\end{center}}{}

\QQTall{会場に座席セットPと座席セットQを作ったところ、机を34台、いすを57脚使った。セットPとセットQをそれぞれ何組作ったか。
\par\smallskip
\begin{center}
\begin{tabular}{|c|c|c|}\hline
 & セットP & セットQ \\ \hline
1組あたりの机 & $2$台 & $4$台 \\ \hline
1組あたりのいす & $6$脚 & $3$脚 \\ \hline
\end{tabular}
\end{center}}{}

\LessonPageBreak
\SetRowSpace{5mm}

